% Work-MAT | 2. rocnik | hodina 20
\documentclass[11pt,a4paper]{article}
\input{../../sablony/pripravy/priprava.sty}
\usepackage{tikz}
\usetikzlibrary{arrows.meta}
\nastavPripravu{2. ročník}{Exponenciální funkce a rovnice}{20}{Průběžný test T2 a rozbor}
\begin{document}
\poznamka{Metodická příprava, NIKOLI hotové zadání testu. Neobsahuje testové varianty ani body. Připrav je v projektové testové šabloně až po schválení skladby a obtížnosti. Orientačně 3 min instrukce, 15--20 min test, 12 min rozbor analogických úloh, 8 min upevnění, zbytek organizace.}
\blok{Před hodinou: příprava ověřování}
\begin{itemize}[label=$\square$]
\item Ověřit, že byly probrány hodiny 11--19; nezkoušet logaritmy ani exponenciální nerovnice, které přijdou později.
\item Připravit varianty, řešení a výtisky podle platných pravidel. Bodové hodnoty a kritéria klasifikace neodvozovat odhadováním.
\item Zkontrolovat aktuální pravidla v Moodle, pomůcky, oznámení a postup pro nepřítomné.
\end{itemize}
\blok{Obsahový rámec testu T2}
Krátký test může ověřit základní vlastnosti funkce $a^x$, čtení grafů včetně asymptot, sestavení předpisu z daných údajů a přepis mocnin na společný základ a volbu jednoduché metody exponenciální rovnice. Podle času lze vybrat více méně náročných úloh nebo několik úloh s postupem; není nutné testovat všechny metody jednou písemkou.
\poznamka{T2 není kopie dřívějších testů z adresáře \texttt{testy/ExponencialniRovnice}. Existující testové soubory je třeba samostatně posoudit, než se přidělí ke konkrétní hodině. Nepřipojuj k této přípravě neověřené úlohy jako hotový test.}
\poznamka{Jde o návrh obsahového okruhu po přesunu témat: hodiny 11--15 jsou výhradně o exponenciálních funkcích, hodiny 16--19 o prvních exponenciálních rovnicích. Na T2 lze klást hlavní důraz na funkce a grafy, jak předpokládal původní tematický plán.}
\blok{Rozbor po odevzdání}
Pro rozbor jsou určeny \emph{analogické nehodnocené} úlohy níže. Neodhaluj při dopisování řešení skutečných variant, pokud budou použity pro náhradní termíny.
\begin{prikladReseniBox}[title={Příklad 1}]
Urči vlastnosti $f(x)=(\frac{1}{4})^x+2$: monotónnost, definiční obor, obor hodnot a asymptotu.
\tcblower
Základ je v $(0;1)$, proto funkce klesá. $D(f)=\mathbb R$, $H(f)=(2;+\infty)$, asymptota $y=2$.
\end{prikladReseniBox}
\begin{prikladReseniBox}[title={Příklad 2}]
Sestav model $g(x)=b\cdot a^x$, jestliže $g(0)=6$ a $g(1)=3$.
\tcblower
$b=6$; ze vztahu $6a=3$ plyne $a=\frac{1}{2}$. Tedy $g(x)=6\cdot\left(\frac{1}{2}\right)^x$.
\end{prikladReseniBox}
\begin{prikladReseniBox}[title={Příklad 3}]
Vyřeš $8^x=4^{x+1}$.
\tcblower
$2^{3x}=2^{2x+2}$, tedy $3x=2x+2$ a $x=2$. Zkouška: $8^2=64$ a $4^3=64$.
\end{prikladReseniBox}
\Needspace{10\baselineskip}
\begin{prikladReseniBox}[title={Příklad 4}]
Vyřeš vytýkáním $3^x+3^{x+1}=108$.
\tcblower
\[
\begin{aligned}
3^x(1+3)&=108,\\
3^x&=27=3^3,\\
x&=3.
\end{aligned}
\]
\end{prikladReseniBox}
\poznamka{Všechna početní řešení lze získat známými mocninami, bez logaritmů nebo substituce. Substituci probíráme od hodiny 21.}
\begin{zadaniReseniBox}[title={Samostatná práce 1}]
\begin{enumerate}[label=\alph*)]
\item $3^x+3^{x+1}=108$;
\item $2^{x+2}+2^x=20$;
\item Kde protne graf $g(x)=2^{x-2}-1$ osu $x$?
\end{enumerate}
\tcblower
\begin{enumerate}[label=\alph*)]
\item $4\cdot3^x=108$, tedy $x=3$.
\item $5\cdot2^x=20$, $2^x=4$, tedy $x=2$.
\item $2^{x-2}=1=2^0$, tedy $x=2$.
\end{enumerate}
\end{zadaniReseniBox}
\begin{zadaniReseniBox}[title={Domácí úkol}]
Pět nehodnocených úloh k upevnění již probraných postupů:
\begin{enumerate}[label=\alph*)]
\item Urči definiční obor a asymptotu $f(x)=2^{x-1}+3$.
\item Urči obor hodnot a monotónnost $g(x)=\left(\frac{1}{3}\right)^x-2$.
\item Vyřeš $8^x=4^{x+1}$.
\item Vyřeš vytýkáním $2^{x+1}+2^x=24$.
\item Vysvětli, proč rovnice $3^x=0$ nemá řešení.
\end{enumerate}
\tcblower
\begin{enumerate}[label=\alph*)]
\item $D(f)=\mathbb R$, vodorovná asymptota $y=3$.
\item $H(g)=(-2;\infty)$, funkce klesá.
\item $2^{3x}=2^{2x+2}$, tedy $x=2$.
\item $3\cdot2^x=24$, tedy $x=3$.
\item Každá mocnina $3^x$ je kladná.
\end{enumerate}
\end{zadaniReseniBox}
\end{document}
